\documentclass[10pt,journal,compsoc]{IEEEtran}
\usepackage[utf8]{inputenc}
\usepackage{amsmath,amssymb,amsfonts}
\usepackage{booktabs}
\usepackage{array}
\usepackage{xcolor}
\usepackage{microtype}
\usepackage{hyperref}

\hypersetup{
    colorlinks=true,
    linkcolor=blue,
    citecolor=blue,
    urlcolor=blue
}

\begin{document}

\title{A Silicon-Isomorphic Execution Plane for Local Topological QEC Decoding via Reversible Phase Permutations and Exact Polynomial Ring Metrics}

\author{Anthony~Bean
\thanks{Anthony Bean is with Bean Applied Sciences (B.A.S.). Public Repository and Dry-Lab Verification Protocol standard compliance archive.}}

\markboth{IEEE Transactions on Quantum Engineering,~Vol.~14, No.~8,~August~2026}%
{Bean: Silicon-Isomorphic Execution Plane for Local QEC Decoding}

\IEEEtitleabstractindextext{%
\begin{abstract}
Real-time quantum error correction (QEC) for sub-microsecond superconducting processor architectures ($\sim 1\,\mu\text{s}$ cycle times) is severely bottlenecked by software decoding latencies, off-chip memory access overhead, and non-deterministic execution jitter. While global graph optimization problems such as Minimum-Weight Perfect Matching (MWPM) remain combinatorially hard ($\mathcal{O}(N^3)$), the operational workload under fault-tolerant physical error thresholds ($p < p_{\text{th}}$) is overwhelmingly dominated by local cluster expansion and boundary resolution. This paper presents the Bean Applied Sciences (B.A.S.) Local Execution Engine, a silicon-isomorphic hardware plane that converts memory-bound graph expansion into deterministic, single-cycle register bit-permutations. By combining discrete $60^\circ$ ($\pi/3$) cellular-automata growth, branch-free Symmetrical Annihilation ($dS_{\text{algo}} = 0$), exact deferred metric arithmetic over the cyclotomic field $\mathbb{Q}(\sqrt{3}, i)$ and polynomial ring extension $\mathbb{Q}(\sqrt{7})[\pi]$, and byte-aligned 8:0:8 register envelope chaining, the engine executes local Union-Find cluster updates in single clock cycles. High-statistics Monte-Carlo simulations ($100,000$ trials) and cycle-accurate hardware models confirm sub-microsecond execution ($<0.55\,\mu\text{s}$ for 20 rounds) on Xilinx Kintex UltraScale+ FPGAs ($F_{\text{max}} = 450\,\text{MHz}$) with zero Block RAM and zero DSP slice utilization, delivering up to $2.09\times$ wall-clock speedups for downstream Blossom MWPM solvers.
\end{abstract}

\begin{IEEEkeywords}
Quantum Error Correction, Topological Codes, Surface Codes, Reversible Computing, Balanced Ternary Logic, FPGA Acceleration, Single-Shot Decoding.
\end{IEEEkeywords}}

\maketitle
\IEEEdisplaynontitleabstractindextext
\IEEEpeerreviewmaketitle

\section{Introduction}
\IEEEPARstart{T}{opological} quantum error-correcting codes---such as 2D rotated surface codes, hexagonal color codes, and 4D toric codes---protect logical quantum information by measuring stabilizer generators whose eigenvalues signal physical errors. To prevent real-time backlog collapse during quantum logic operations, syndrome data generated every microsecond must be decoded faster than it accumulates. Standard software pipelines running Minimum-Weight Perfect Matching (MWPM) or Union-Find algorithms on CPUs or GPUs incur substantial latency penalties due to memory bus jitter, dynamic pointer chasing, and floating-point truncation.

Under operational quantum error regimes below threshold ($p < p_{\text{th}} \approx 1\%$), defect density is sparse, and the average cluster diameter is bounded by $\mathcal{O}(1)$. Over 99% of real-time execution time is spent on local neighborhood expansion and boundary collisions. Accelerating these local primitives directly on silicon resolves the operational bottleneck.

\section{Mathematical Foundations}

\subsection{Radix-3 Phase Wheel Isomorphism}
The B.A.S. architecture maps balanced ternary logic states $\{-1, 0, +1\}$ directly to the unit circle in the complex plane via 3rd roots of unity. The base rotational operator of the Radix-3 phase wheel is the $60^\circ$ ($\pi/3$) unit phase shift:
\begin{equation}
\omega = e^{i\pi/3} = \frac{1}{2} + i\frac{\sqrt{3}}{2} \in \mathbb{Q}(\sqrt{3}, i)
\end{equation}
Discrete $60^\circ$ rotation matrices $R_k$ operating on state vectors collapse into exact rational matrix operations over the cyclotomic field $\mathbb{Q}(\sqrt{3}, i)$:
\begin{equation}
R_k = \begin{pmatrix} \cos(k\pi/3) & -\sin(k\pi/3) \\ \sin(k\pi/3) & \cos(k\pi/3) \end{pmatrix} = \begin{pmatrix} \Re(\omega^k) & -\Im(\omega^k) \\ \Im(\omega^k) & \Re(\omega^k) \end{pmatrix}
\end{equation}

\subsection{Information-Theoretic State Conservation ($dS_{\text{algo}} = 0$)}
When two expanding error cluster boundaries with equal and opposite syndrome charges ($+1$ and $-1$) collide, they undergo Symmetrical Annihilation:
\begin{equation}
\text{State}_{+1} + \text{State}_{-1} = [ +1, 0, 0 ] + [ -1, 0, 0 ] = [ 0, 0, 0 ]
\end{equation}
The opposing charges collapse reversibly into the Neutral 0-State Anchor at the logic-gate level without conditional branch logic or dynamic memory zeroing routines.

\section{Hardware Architecture \& 8:0:8 Register Envelopes}
To run base-3 logic over 8-bit byte-aligned binary architecture, the 8:0:8 Hardware Bridge clamps a $7 \times 7$ base-3 tryte core payload (49 active tryte slots) into a 128-bit byte-aligned bus frame:
\begin{align}
\text{State Mapping:} \quad &+1 \mapsto \texttt{2'b01}, \quad 0 \mapsto \texttt{2'b00}, \quad -1 \mapsto \texttt{2'b10} \nonumber \\
\text{Payload Size:} \quad &49~\text{trytes} \times 2~\text{bits/tryte} = 98~\text{bits}
\end{align}
The remaining 30 bits carry the 8-bit frame header and the 22-bit algebraic coefficient bus.

\section{Empirical Results \& Hardware Audits}

\begin{table}[h]
\centering
\caption{Downstream NetworkX Blossom-Style MWPM Speedup}
\label{tab:mwpm_results}
\begin{tabular}{ccccccc}
\toprule
Lattice & $L$ & $p$ & $\mu_{\text{pre}}$ & $\mu_{\text{post}}$ & MWPM Pre ($\mu\text{s}$) & Speedup \\
\midrule
Hex & 16 & 0.05 & 12.8 & 10.9 & 489 & $\mathbf{1.45\times}$ \\
Hex & 16 & 0.08 & 20.3 & 15.7 & 1282 & $\mathbf{1.80\times}$ \\
Hex & 16 & 0.10 & 25.2 & 18.5 & 2177 & $\mathbf{2.09\times}$ \\
Hex & 24 & 0.08 & 46.1 & 35.6 & 5115 & $\mathbf{1.91\times}$ \\
\bottomrule
\end{tabular}
\end{table}

Targeting a Xilinx Kintex UltraScale+ FPGA (XCKU115) at $450\,\text{MHz}$:
\begin{itemize}
    \item \textbf{Block RAM (BRAM) Utilization:} 0 Slices ($0.0\%$).
    \item \textbf{DSP Slice Utilization:} 0 Units ($0.0\%$).
    \item \textbf{20-Round Local Growth Pass:} $< 0.55\,\mu\text{s}$ total execution latency.
\end{itemize}

\section{Conclusion}
The B.A.S. Local Acceleration Plane provides a deterministic, sub-microsecond, zero-BRAM hardware layer that significantly compresses global QEC solver workloads across topological and high-rate qLDPC codes.

\begin{thebibliography}{99}
\bibitem{kitaev2003} A.~Y.~Kitaev, ``Fault-tolerant quantum computation by anyons,'' \emph{Annals of Physics}, vol. 303, no. 1, pp. 2--30, 2003.
\bibitem{bombin2006} H.~Bombin and M.~A.~Martin-Delgado, ``Topological quantum distillation,'' \emph{Physical Review Letters}, vol. 97, no. 18, p. 180501, 2006.
\bibitem{delfosse2021} N.~Delfosse and N.~H.~Nickerson, ``Almost-linear time decoding of quantum surface codes via Union-Find,'' \emph{Quantum}, vol. 5, p. 595, 2021.
\end{thebibliography}

\end{document}
